% Generated from selected worked problems in committed QFT-soln sources.
% Source repository: https://github.com/Luca-Yucheng-Jin/QFT-soln
% Source commit: d74e71bc8fb585141e81696c5d2b810008d9dba6
\section{Wald General Relativity, Chapter 6: The Schwarzschild Solution}
\markboth{\thesection\quad WALD GENERAL RELATIVITY, CHAPTER 6}{}

\noindent\textit{Problem prompts are paraphrased from Robert M. Wald,
\emph{General Relativity} (University of Chicago Press, 1984); the equations,
notation, numbering, and guidance follow the source. See the
\href{https://press.uchicago.edu/ucp/books/book/chicago/G/bo5952261.html}{official publisher page}.}

\subsection{Chapter 6, Problem 1: Isotropic Coordinates}

Let $M$ be a three-dimensional manifold with a spherically symmetric Riemannian
metric, and suppose $\nabla_a r\neq0$, where $r$ is defined in (6.1.3).

\begin{enumerate}[label=\alph*)]
    \item Show that one may introduce an isotropic radial coordinate $\widetilde r$
    for which
    \begin{equation}
        ds^2=H(\widetilde r)
        \left[d\widetilde r^2+\widetilde r^2d\Omega^2\right].
    \end{equation}
    Thus every spherically symmetric three-dimensional space is conformally flat.
    \begin{solution}
        Assuming the metric can be written in a conformally flat form, then we must have
        \begin{equation}
            h(r)dr^2  = H(\tilde r)d\tilde r^2, \quad H(\tilde r)\tilde r^2 = r^2,
        \end{equation}
        which gives the differential equation
        \begin{equation}
            \frac{\sqrt{h(r)}}{r}dr = \frac{d \tilde r}{\tilde r} \implies \tilde r = A \exp\left( \int^r dr' \frac{\sqrt{h(r')}}{r'}\right).
        \end{equation}

    \end{solution}

    \item Express the Schwarzschild metric in these coordinates and obtain
    \begin{equation}
        ds^2
        =-\frac{(1-M/2\widetilde r)^2}{(1+M/2\widetilde r)^2}\,dt^2
        +\left(1+\frac{M}{2\widetilde r}\right)^4
        \left[d\widetilde r^2+\widetilde r^2d\Omega^2\right].
    \end{equation}
    \begin{solution}
        For the Schwarzschild metric,
        \begin{equation}
            \tilde r = A \exp\left(\int ^r dr' \frac{1}{\sqrt{r'^2 - 2Gmr'}}\right) = 2A\left(r-GM + \sqrt{r(r-GM)}\right).
        \end{equation}
        $A$ can be fixed by imposing the metric to be asymptotically flat, that is as $r \to \infty$, $H(\tilde r ) \to 1$. Therefore, we must have $A = 1/4$.
        Hence, playing around with the algebra slightly, we find
        \begin{equation}
            \tilde r = \tilde r\left(1 + \frac{GM}{2\tilde r}\right)^2,
        \end{equation}
        which gives
            \begin{equation}
        ds^2
        =-\frac{(1-M/2\widetilde r)^2}{(1+M/2\widetilde r)^2}\,dt^2
        +\left(1+\frac{M}{2\widetilde r}\right)^4
        \left[d\widetilde r^2+\widetilde r^2d\Omega^2\right].
    \end{equation}
    \end{solution}
\end{enumerate}

\subsection{Chapter 6, Problem 3: The Reissner--Nordstr\"om Solution}

Consider the source-free Maxwell equations (4.3.12) and (4.3.13), with $j^a=0$,
on the static, spherically symmetric spacetime (6.1.5).

\begin{enumerate}[label=\alph*)]
    \item Argue that a Maxwell tensor with the same static and spherical
    symmetries must have the form
    \begin{equation}
        F_{ab}
        =2A(r)(e_0)_{[a}(e_1)_{b]}
        +2B(r)(e_2)_{[a}(e_3)_{b]},
    \end{equation}
    where the $(e_\mu)_a$ are defined in (6.1.6).
    \begin{solution}
        Since $F$ is a two form, it should in general take the form
        \begin{equation}
            F = F_{ab} \hat \theta^a \wedge \hat \theta^b,
        \end{equation}
        where $\hat \theta^a$ is the dual basis of $\hat e_a$. Now, excluding the basis that are not invariant under $SO(3)$ transformation, we are left with
        \begin{equation}
            F = A(r) \hat \theta^0 \wedge \hat \theta^1 + B(r) \hat{\theta}^2 \wedge\hat \theta^3,
        \end{equation}
        where $A, B$ are only functions of $r$ by the requirement of the universe being static.
    \end{solution}

    \item Restrict to $B(r)=0$ and solve Maxwell's equations to find
    \begin{equation}
        A(r)=-\frac{q}{r^2},
    \end{equation}
    interpreting $q$ as the total charge. The branch with $B(r)\neq0$ is a
    duality rotation of this electric solution and represents a magnetic monopole.
    \begin{solution}
        Writing the field strength two form in terms of the coordinate basis, we have
        \begin{equation}
            F = A(r) f^{1/2}h^{1/2}dt\wedge dr.
        \end{equation}\
        Since $A, f ,h$ are all functions of $r$, $dF = 0$ gives no extra informations. The wedge of $F$ is given by
        \begin{equation}
            \star F = A(r)r^2\sin\theta d\theta \wedge d\phi,
        \end{equation}
        up to some constant that would not affect our calculation here.
        Then, the source-free Maxwell equation gives
        \begin{equation}
            \partial_r(A(r)r^2) = 0 \implies A(r) = -\frac{q}{r^2}.
        \end{equation}

    \end{solution}

    \item Insert the electromagnetic stress-energy tensor from part (b) into
    \begin{equation}
        G_{ab}=8\pi T_{ab},
    \end{equation}
    solve the resulting equations, and show that the general metric is
    \begin{equation}
        ds^2
        =-\left(1-\frac{2M}{r}+\frac{q^2}{r^2}\right)dt^2
        +\left(1-\frac{2M}{r}+\frac{q^2}{r^2}\right)^{-1}dr^2
        +r^2d\Omega^2.
    \end{equation}
    \begin{solution}
        The resulting energy momentum tensor in vierbein basis is
        \begin{equation}
            T_{00} = \frac{1}{8\pi}\frac{q^2}{r^4},
        \end{equation}
        \begin{equation}
            T_{11} = -\frac{1}{8\pi}\frac{q^2}{r^4},
        \end{equation}
        \begin{equation}
            T_{22} = \frac{1}{8\pi}\frac{q^2}{r^4},
        \end{equation}
        and
        \begin{equation}
            T_{33} = \frac{1}{8\pi}\frac{q^2}{r^4}.
        \end{equation}
        The Einstein equation is given by
        \begin{equation}
            R_{00} +\frac{1}{2}R = 8\pi T_{00},
            \label{eq:00}
        \end{equation}
        \begin{equation}
            R_{11} - \frac{1}{2}R = 8\pi T_{11},
            \label{eq:11}
        \end{equation}
        \begin{equation}
            R_{22} - \frac{1}{2}R = 8\pi  T_{22},
            \label{eq:22}
        \end{equation}
        and
        \begin{equation}
            R_{33} - \frac{1}{2}R = 8\pi T_{33},
            \label{eq:33}
        \end{equation}
        where the Ricci tensor is given in (6.1.35-6.1.37).
        Combing \eqref{eq:00} and \eqref{eq:11} , we find
        \begin{equation}
            R_{00} + R_{11} = 0,
        \end{equation}
        which is equivalent to (6.1.38), and hence
        \begin{equation}
            f = h^{-1}.
        \end{equation}
        Note that \eqref{eq:22} and \eqref{eq:33} are completely identical, plugging $f=1/h$ in one of them, we find
        \begin{equation}
            f''  + \frac{2f'}{r}= \frac{2q^2}{r^4},
        \end{equation}
        where we have used
        \begin{equation}
            R = -R_{00} + R_{11} + 2R_{22}.
        \end{equation}
        This gives
        \begin{equation}
            f = h^{-1} = 1- \frac{C}{r} -\frac{q^2}{r^2}.
        \end{equation}


    \end{solution}
\end{enumerate}

\subsection{Chapter 6, Problem 4: Local Force and Redshift}

Let $(M,g_{ab})$ be stationary with timelike Killing field $\xi^a$, and define
\begin{equation}
    V^2=-\xi^a\xi_a.
\end{equation}

\begin{enumerate}[label=\alph*)]
    \item For a stationary observer, show that the acceleration
    $a^b=u^a\nabla_a u^b$ is
    \begin{equation}
        a^b=\nabla^b\ln V.
    \end{equation}
    \begin{solution}
        \begin{equation}
            a^b = u^a \nabla_a u^b  = \frac{u^a \nabla_a \xi^b}{V} - \frac{u^a \xi^b \nabla_a V}{V^2}.
        \end{equation}
        The second term vanishes by
        \begin{equation}
            \frac{u^a \xi^b \nabla_a V}{V^2} = -\frac{2\xi^a\xi^b\xi^c\nabla_a\xi^c}{V^3} = 0
        \end{equation}
        by $\nabla_a \xi_b = - \nabla_b \xi_a$. Then, we have
        \begin{equation}
            a^b = \frac{\xi^a \nabla_a \xi^b}{V^2} = -\frac{\xi^a\nabla^b \xi_a}{V^2} = \frac{1}{2}\nabla^b \log V^2 = \nabla^b \log V.
        \end{equation}
    \end{solution}

    \item Suppose additionally that the spacetime is asymptotically flat. Thus
    there are coordinates $t,x,y,z$, with
    $\xi^a=(\partial/\partial t)^a$, in which $g_{ab}$ tends to
    $\operatorname{diag}(-1,1,1,1)$ as
    $r=(x^2+y^2+z^2)^{1/2}\to\infty$.
    See Chapter 11 for further discussion of asymptotic flatness.

    A particle of mass $m$ has energy at infinity
    \begin{equation}
        E=-m\xi^a u_a.
    \end{equation}
    Hold the particle stationary with a massless string whose far end is held by
    a stationary observer at large $r$. If $F$ is the force applied locally to the
    particle, part (a) gives
    \begin{equation}
        F=mV^{-1}\left(\nabla_aV\nabla^aV\right)^{1/2}.
    \end{equation}
    Use conservation of energy to show that the force supplied at infinity is
    \begin{equation}
        F_\infty=VF.
    \end{equation}
    Interpret the factor $V$ as the redshift relating the two force magnitudes.
    \begin{solution}
        The local magnitude of force is given by
        \begin{equation}
            F = m \sqrt{a^a a_a} = mV^{-1}\left(\nabla_aV\nabla^aV\right)^{1/2},
        \end{equation}
        while the energy of a stationary particle measured by a local stationary observer is given by
        \begin{equation}
            E = m.
        \end{equation}
        Assuming the stationary particle is pushed through a proper distance $\delta s$ adiabatically with a constant force $F$. Then, the work done locally is equal to the local change of energy is given by
        \begin{equation}
            F\delta s = \delta E.
        \end{equation}
        For an observer at spatial infinity, he agrees on the proper distance, while the force seen by this observer, $F_\infty$, is different. Therefore, the change of energy seen by this observer is
        \begin{equation}
            \delta E_\infty = F_\infty \delta s.
        \end{equation}
        By the energy redshift, we have
        \begin{equation}
            V\delta E = \delta E_\infty.
        \end{equation}
        Therefore, we have
        \begin{equation}
            F_\infty =VF.
        \end{equation}
    \end{solution}
\end{enumerate}

\subsection{Chapter 6, Problem 5: Relativistic Time Delay}

Derive the general-relativistic time-delay formula (6.3.45).
\emph{Draft solution: the derivation below does not yet display the final time-delay formula (6.3.45).}
\begin{solution}
    By time translational symmetry, the time delay from the Earth to the planet in coordinate time must be equivalent to
    \begin{equation}
        \Delta t = \int^{R_E}_{R_0} dr\, \frac{1}{(1-2GM/r)\sqrt{1-(1-2GM/r)b^2/r^2}} + \int^{R_p}_{R_0} dr\, \frac{1}{(1-2GM/r)\sqrt{1-(1-2GM/r)b^2/r^2}}.
    \end{equation}
    To first order in $G$, this is
    \begin{equation}
        \Delta t = \Delta t \bigg|_{G=0} + \frac{\partial (\Delta t)}{\partial G}\bigg|_{G=0}G + \mathcal{O}(G^2).
    \end{equation}
    The first term is simply the time elapsed when the star is absent. Therefore, it must be
    \begin{equation}
        \Delta t \bigg|_{G=0} = (R_E^2-R_0^2)^{1/2} + (R_p^2 - R_0^2)^{1/2}
    \end{equation}
    by simple geometry. Recall that
    \begin{equation}
        R_0^3 - b^2(R_0-2GM) = 0 \implies b = R_0 + GM + \mathcal{O}(G^2).
    \end{equation}
    Then, for the second term, we find
    \begin{equation}
    \begin{aligned}
        \frac{\partial (\Delta t)}{\partial G}\bigg|_{G=0}G &= \frac{2GM}{R_0}\log\left(\frac{r}{R_0} + \frac{\sqrt{r^2-R_0^2}}{R_0}\right)\bigg|^{R_E}_{R_0}+\frac{2GM}{R_0}\log\left(\frac{r}{R_0} + \frac{\sqrt{r^2-R_0^2}}{R_0}\right)\bigg|^{R_p}_{R_0}\\
        \qquad &-\frac{2GMr}{\sqrt{r^2 -b^2}}\bigg|^{R_E}_{R_0}-\frac{2GMr}{\sqrt{r^2 -b^2}}\bigg|^{R_p}_{R_0}.
    \end{aligned}
    \end{equation}
    Note that special attention is required when substituting $R_0$ into the last two terms. Carefully keeping terms of first order in $G$, one would eventually find the residue terms of (6.3.45).
\end{solution}

\subsection{Chapter 6, Problem 6: Lifetime Inside the Schwarzschild Horizon}

Show that every particle in region II, $r<2M$, of the extended Schwarzschild
spacetime in Figure 6.9, whether geodesic or accelerated, must obey
\begin{equation}
    \left|\frac{dr}{d\tau}\right|
    \geq\left(\frac{2M}{r}-1\right)^{1/2}.
\end{equation}
Deduce that no observer in region II can have a remaining proper lifetime greater
than
    \begin{equation}
        \tau=\pi M
        \sim 10^{-5}\left(\frac{M}{M_\odot}\right)\mathrm{s}
    \end{equation}
before reaching $r=0$. Finally, show that this bound is approached by free fall
from $r=2M$ in the limit $E\to0$.
\begin{solution}
    Directly from the Schwarzschild metric, we immediately
    \begin{equation}
        \left|\frac{dr}{d\tau}\right| \geq \left|1-\frac{2GM}{r}\right|^{1/2} = \left(\frac{2GM}{r}- 1\right)^{1/2}
    \end{equation}
    for $r < 2GM$. This gives
    \begin{equation}
        \tau \leq \int_0^{2GM}\frac{dr}{\sqrt{2GM/r-1}} = \pi GM.
    \end{equation}
    The integral is precisely the proper time for a particle to fall to the singularity, released from the event horizon with $ E = 0$.
\end{solution}
